% NeurIPS 2026 — course submission (preprint mode)
\documentclass{article}

% NeurIPS 2026 style (file must be in same directory at compile time)
\usepackage[preprint]{neurips_2026}

% Standard encoding and font
\usepackage[utf8]{inputenc}
\usepackage[T1]{fontenc}

% Hyperlinks
\usepackage{hyperref}
\usepackage{url}

% Tables
\usepackage{booktabs}

% Math
\usepackage{amsfonts}
\usepackage{amsmath}
\usepackage{nicefrac}
\usepackage{microtype}

% Figures
\usepackage{graphicx}

% Color
\usepackage{xcolor}

% Citations
\usepackage{natbib}

% -----------------------------------------------------------------------
\title{Automated Feature Discovery for National Park Biodiversity\\
       Prediction: An AutoResearch Study}

\author{%
  Isabella Woods \\
  STAT 390 Capstone, Northwestern University \\
  \texttt{isabellawoods2026@u.northwestern.edu}
}

\begin{document}

\maketitle

% -----------------------------------------------------------------------
\begin{abstract}
We apply an AutoResearch loop---a human-supervised agent iteration
protocol---to the problem of predicting per-taxon Shannon Diversity
Index~(SDI) in 149 US National Parks using iNaturalist citizen-science
observations. Starting from a Ridge regression baseline
(RMSE~$=$~0.326), 85 disciplined experiments reduced validation RMSE
to 0.238 ($R^{2}$~$=$~0.778), a 27\% improvement. The most impactful
interventions were autoregressive temporal lags, cumulative SDI
baselines, and taxon-specific ecological interaction features
discovered via greedy forward selection. We demonstrate that linear
regularized models generalize better than tree-based models under
park-based cross-validation, and that the remaining performance gap
reflects irreducible cross-park heterogeneity rather than model
capacity.
\end{abstract}

% -----------------------------------------------------------------------
\section{Introduction}
\label{sec:intro}

National Park Service (NPS) biodiversity monitoring traditionally
relies on structured field surveys that are expensive, infrequent, and
geographically uneven. Citizen-science platforms such as iNaturalist
now provide millions of georeferenced observations annually, but raw
observation counts are severely confounded by observer effort: the
platform grew more than 20$\times$ between 2010 and 2025, producing an
apparent upward trend in apparent diversity ($r = 0.88$ between
$\log(\text{observations})$ and SDI) that is methodological rather than
ecological.

We address this problem in two ways. First, rarefaction---subsampling
each (park, month, taxon) cell to a fixed $N = 50$ observations before
computing SDI---removes the effort bias and reduces the spurious
correlation to $r = 0.35$. Second, we apply an \emph{AutoResearch
loop}: a structured agent-assisted iteration protocol in which a large
language model proposes feature engineering hypotheses, edits a
designated section of the training script, evaluates the result, and
commits to a version-controlled experiment log. Human oversight is
retained for problem framing, leakage detection, and stopping
decisions.

\textbf{Contributions.} (1) We achieve a 27\% validation RMSE
reduction over 85 experiments (0.326 $\to$ 0.238), with autoregressive
temporal lags contributing the largest single gain
($\Delta\text{RMSE} = 0.073$). (2) We provide an empirical
demonstration that regularized linear models outperform gradient
boosting under strict park-based cross-validation on this dataset.
(3) We characterize the failure modes of the AutoResearch protocol,
including leakage architecture errors and diminishing-returns pursuit
past ecological significance.

\textbf{Research question.} Can an LLM-assisted AutoResearch loop
discover effective feature engineering strategies for time-series
biodiversity prediction, starting only from a minimal baseline?

% -----------------------------------------------------------------------
\section{Data and Problem Formulation}
\label{sec:data}

\subsection{Dataset}

The raw dataset consists of iNaturalist Research-Grade observations
\citep{inaturalist2025} spanning 2021--2025 for 149 US National Parks.
Observations are aggregated by (park, year, month, taxon group), where
taxon group is one of 13 categories: Aves, Mammalia, Plantae, Insecta,
Reptilia, Amphibia, Fungi, Arachnida, Mollusca, Actinopterygii,
Animalia, and two residual groups. These are joined with MODIS-derived
satellite features (fraction of photosynthetically active radiation
[FPAR], gross primary production [GPP], evapotranspiration [ET], snow
cover [SNOW]) and park-level human-activity features (annual visitor
counts, traffic volume, population density within a 50\,km radius, and
facility counts).

\subsection{Rarefaction}

Each (park, month, taxon) cell is subsampled to exactly $N = 50$
observations before computing SDI, using a fixed random seed for
reproducibility. Cells with fewer than 50 observations are excluded.
This rarefaction step \citep{hurlbert1971} reduces the effort
correlation from $r = 0.88$ to $r = 0.35$ and eliminates the
spurious temporal trend visible in raw SDI. The resulting dataset
contains \textbf{8,824 rows}, each a (park, year, month, taxon)
quadruple.

The target variable is the rarefied Shannon Diversity Index:
\begin{equation}
  \text{SDI}_{\text{rarefied}} = -\sum_{i} p_{i} \log p_{i},
  \label{eq:sdi}
\end{equation}
where $p_i$ is the proportion of species $i$ within the rarefied
subsample.

\subsection{Train/Validation Split}

We use a \textbf{park-based} split: all months and taxa from a given
park are assigned entirely to either the training set (108 parks) or
the validation set (27 parks). This prevents any temporal or taxonomic
leakage across parks and mirrors the intended use case, where
predictions must generalize to parks not seen during model fitting. The
evaluation metric throughout is validation RMSE on the 27 held-out
parks (never used during feature engineering or model selection).

% -----------------------------------------------------------------------
\section{AutoResearch Loop Design}
\label{sec:loop}

\subsection{Protocol}

Each experiment follows a fixed protocol: (1) read the current state
of \texttt{train.py} and \texttt{results.tsv}; (2) form a hypothesis
about a feature or model change likely to reduce validation RMSE; (3)
edit only the designated feature-engineering section (\S2) of
\texttt{train.py}; (4) run the script with a 10-minute kill timer; (5)
compare validation RMSE to the previous best; (6) keep the change if
RMSE strictly improves, otherwise revert; (7) append the result to
\texttt{results.tsv} and commit to Git with a descriptive message.
The evaluation code (train/val split, metric computation, random seeds)
is never modified. This append-only log and single-change discipline
prevent compounding confounds and enable clean ablations.

\subsection{Feature Engineering Space}

Over 85 experiments, the agent explored the following feature families:

\begin{itemize}
  \item \textbf{ECO} --- satellite-derived ecological state: FPAR,
    GPP, ET, SNOW, fire indicators.
  \item \textbf{HUM} --- human-impact proxies: visitor counts,
    traffic, population density, facility counts.
  \item \textbf{GEO} --- geographic position: latitude, longitude.
  \item \textbf{TIME} --- temporal encoding: $\sin/\cos$ of month
    (seasonality), \texttt{year\_norm} (linear trend).
  \item \textbf{AR lags} --- autoregressive features: SDI at lag
    1, 2, 3, 6, 9, and 12 months within the same (park, taxon) series.
  \item \textbf{Cumulative baselines} --- expanding mean SDI by park
    $\times$ taxon (\texttt{SDI\_cumean}), expanding mean of the same
    calendar month (\texttt{SDI\_month\_cumean}), and park-level
    cross-taxon mean (\texttt{SDI\_park\_cumean}).
  \item \textbf{Deviations} --- difference between lag-1 and seasonal
    norm (\texttt{SDI\_dev}), interannual deviation (\texttt{SDI\_dev12}),
    and prior-observation reliability count (\texttt{n\_prior\_months}).
  \item \textbf{Taxon interactions} --- products of taxon-group
    indicator with FPAR, latitude, population density, season, and
    year trend, selected by greedy forward search over 30 candidates.
\end{itemize}

Experiments ran in 5 phases: alpha/model tuning (Exp.\ 032--041),
AR lag expansion (Exp.\ 042--059), ElasticNet and baselines
(Exp.\ 060--064), deviation and interaction features (Exp.\ 065--075),
and final validation and ablations (Exp.\ 076--085).

% -----------------------------------------------------------------------
\section{Experiments and Results}
\label{sec:results}

\begin{table}[t]
  \caption{Key experimental phases. Val RMSE is measured on 27
    held-out parks. ``Experiments'' column gives the range of
    experiment indices in the run log.}
  \label{tab:phases}
  \centering
  \small
  \begin{tabular}{llllrr}
    \toprule
    Phase & Exp.\ range & Key change & Model & Val RMSE & Val $R^{2}$ \\
    \midrule
    Baseline       & 032      & Ridge($\alpha$=100), rarefied SDI           & Ridge       & 0.326 & 0.611 \\
    AR lag-1       & 042--044 & SDI\_lag1 prior-month signal                & Ridge/Lasso & 0.278 & 0.722 \\
    AR lag-2 to 12 & 054--059 & SDI\_lag2,3,6,9,12 multi-scale AR           & Lasso       & 0.253 & 0.750 \\
    Baselines      & 063--064 & SDI\_cumean, SDI\_month\_cumean             & ElasticNet  & 0.246 & 0.764 \\
    Deviations     & 066--070 & SDI\_dev, SDI\_dev12, taxon$\times$season   & ElasticNet  & 0.243 & 0.768 \\
    Interactions   & 072--075 & taxon$\times$FPAR/lat/pop + greedy (7/30)   & ElasticNet  & 0.238 & 0.778 \\
    Final          & 079      & SDI\_park\_cumean park-level baseline        & ElasticNet  & 0.238 & 0.778 \\
    \bottomrule
  \end{tabular}
\end{table}

\subsection{Final Model}

The best model is \texttt{ElasticNet(alpha=0.001, l1\_ratio=0.2,
max\_iter=10000)} applied to 195 numeric features plus one
one-hot-encoded taxon categorical variable. The $\ell_1$ ratio of 0.2
(more Ridge-like than Lasso) is consistent with the high mutual
correlation among AR lags, cumulative baselines, and deviation
features: mild sparsity pressure reduces coefficients but does not
zero them. Stability is confirmed by two independent runs with
\texttt{random\_state=42} yielding identical validation RMSE to 6
decimal places.

\subsection{Metric Trajectory}

Figure~\ref{fig:trajectory} shows the full validation RMSE and $R^{2}$
trajectory across all 85 experiments. The AR lag phase (Exp.\ 042--059)
produces the steepest decline, while subsequent phases yield diminishing
but cumulatively meaningful returns.

\begin{figure}[h]
  \centering
  \includegraphics[width=0.85\linewidth]{outputs/metric_over_time.png}
  \caption{Validation RMSE (left axis, blue) and $R^{2}$ (right axis,
    orange) across all 85 experiments. Filled circles indicate kept
    experiments; open circles indicate discarded experiments. The
    dashed horizontal line marks the 0.20 target RMSE. The largest
    single drop corresponds to the introduction of SDI\_lag1
    (Exp.\ 042).}
  \label{fig:trajectory}
\end{figure}

\subsection{Residual Diagnostics}

Figure~\ref{fig:residuals} summarizes residual diagnostics for the
final model. Residuals are approximately centered around zero across
most parks, with elevated variance in parks that have sparse temporal
coverage (fewer than 6 months with $\geq$50 observations per taxon).
No systematic bias by taxon group is observed, suggesting that the
taxon interaction features successfully capture the main differential
responses.

\begin{figure}[h]
  \centering
  \includegraphics[width=0.85\linewidth]{outputs/residual_diagnostics.png}
  \caption{Residual diagnostics for the final ElasticNet model on
    the 27 held-out validation parks. Top row: residual vs.\ fitted
    values and QQ-plot. Bottom row: residuals by taxon group and by
    calendar month, showing no systematic bias in either dimension.}
  \label{fig:residuals}
\end{figure}

\subsection{What Failed}

Several experimental directions were explored and rejected:

\textbf{Tree-based models.}
\texttt{HistGradientBoostingRegressor} consistently produced worse
validation RMSE (0.250--0.303) despite better or comparable CV RMSE
(0.210--0.236). The validation--CV gap ($\Delta \approx 0.04$) indicates
that gradient boosting memorizes park-specific patterns in the training
set that do not transfer to held-out parks. Regularized linear models,
by contrast, cannot represent park-level intercepts and are therefore
forced to generalize.

\textbf{PCA compression.}
Compressing the taxon interaction block with PCA(50) degraded
validation RMSE from 0.238 to 0.240. The interaction features are
already sparse and semantically structured; dimensionality reduction
destroyed this structure.

\textbf{Park $\times$ taxon target encoding.}
Encoding each (park, taxon) pair as its training-set mean SDI caused
catastrophic generalization failure (RMSE = 0.481 on validation). For
the 27 held-out parks, the encoder imputed the global median---a poor
substitute for park-specific context. This is a canonical example of
target-encoding leakage that required human detection to diagnose.

\textbf{Generalization gap.}
The final model achieves CV RMSE = 0.210 vs.\ validation RMSE = 0.238
($\Delta = 0.028$). This modest gap indicates that cross-park
heterogeneity is real but not catastrophic; the park-based CV
procedure is reasonably calibrated.

% -----------------------------------------------------------------------
\section{Discussion and Limitations}
\label{sec:discussion}

\subsection{Remaining Performance Gap}

The final model is 0.038 RMSE above the 0.20 target. We attribute this
primarily to \emph{cross-park heterogeneity}: each park occupies a
distinct position in a high-dimensional ecological space defined by
elevation, climate zone, dominant vegetation biome, and management
history. None of these variables are directly available in our feature
set, and their proxies (latitude, FPAR, snow cover) provide only a
partial representation.

\subsection{Linear Model Requirement}

The consistent superiority of ElasticNet over gradient boosting under
park-based cross-validation reflects a structural property of the
problem: with only 108 training parks, tree-based models have
sufficient capacity to memorize park-specific response functions,
including idiosyncratic seasonal patterns and taxon composition, that
do not appear in the 27 validation parks. Regularized linear models
cannot express such memorization and are therefore forced to rely on
transferable patterns.

\subsection{Agent Limitations}

Three recurring failure modes were identified in the AutoResearch
protocol:

\begin{enumerate}
  \item \textbf{Leakage architecture errors.} The target-encoding
    experiment (Exp.\ 081--081b) introduced a train/test leakage
    that the agent did not independently detect. Human review of the
    result (RMSE = 0.481, more than double the baseline) was required
    to identify and reject the change.
  \item \textbf{Diminishing returns pursuit.} After Exp.\ 075
    (RMSE = 0.238), the agent continued for 10 additional experiments
    with improvements smaller than $10^{-3}$ RMSE. These experiments
    consumed runtime without ecological significance.
  \item \textbf{Greedy selection limitations.} Forward selection over
    30 taxon interaction candidates retained 7. This process cannot
    capture synergistic interactions between features not individually
    informative, potentially missing combined signals.
\end{enumerate}

\subsection{Recommended Redesign}

In retrospect, the most impactful architectural change would have been
to introduce a mixed-effects model structure from the outset, with
park-level random intercepts. This would have separated park-specific
baselines from transferable ecological signals, reducing the tension
between memorization and generalization that constrained the tree-based
models. Climate zone and elevation strata as grouping variables would
further stratify the heterogeneity.

% -----------------------------------------------------------------------
\section{Conclusion}
\label{sec:conclusion}

Starting from a Ridge regression baseline (RMSE = 0.326, $R^{2}$ =
0.611), 85 disciplined AutoResearch experiments reduced validation RMSE
to 0.238 ($R^{2}$ = 0.778), a 27\% improvement. The dominant driver
was autoregressive temporal context: SDI lags at 1, 2, 3, 6, 9, and
12 months, together with expanding-mean baselines, contributed an RMSE
reduction of 0.080. Taxon-specific ecological interaction features
provided an additional 0.008 RMSE reduction via greedy forward
selection.

The study confirms that regularized linear models generalize better
than gradient boosting under strict park-based cross-validation when
the number of parks is small ($n = 108$ training parks). It also
illustrates that agent-guided research is effective for systematic
feature search but requires human oversight for problem framing,
leakage detection, and stopping decisions based on ecological
significance rather than numerical improvement alone.

Experiment logs, training code, and results are maintained on the
project Git repository at branch \texttt{autoresearch/apr27}.

% -----------------------------------------------------------------------
\bibliographystyle{abbrvnat}
\begin{thebibliography}{9}

\bibitem[Shannon(1948)]{shannon1948}
Shannon, C.~E. (1948).
\newblock A mathematical theory of communication.
\newblock \emph{Bell System Technical Journal}, 27(3):379--423.

\bibitem[Hurlbert(1971)]{hurlbert1971}
Hurlbert, S.~H. (1971).
\newblock The nonconcept of species diversity: a critique and alternative
  parameters.
\newblock \emph{Ecology}, 52(4):577--586.

\bibitem[Zou and Hastie(2005)]{zou2005}
Zou, H. and Hastie, T. (2005).
\newblock Regularization and variable selection via the elastic net.
\newblock \emph{Journal of the Royal Statistical Society: Series B},
  67(2):301--320.

\bibitem[iNaturalist(2025)]{inaturalist2025}
{iNaturalist} (2025).
\newblock iNaturalist Research-Grade Observations.
\newblock \url{https://www.inaturalist.org/}.
\newblock Accessed 2025.

\end{thebibliography}

\end{document}
